\section*{Abstract}
FFGFT exists in several mathematical languages -- torus geometry, matrix algebra, Hilbert-space
operators, information formalism, recursion, two layers. This document does not ask how FFGFT
stands relative to \emph{other} frameworks (Doc.~265 does that outward); it asks how these internal
representations stand \emph{among themselves}. Finding: they are not a flat set of equals but a
\emph{generative core} (the $T^4$ torus) plus three clearly distinct relationship types --
\textbf{identity}, \textbf{bijection}, \textbf{projection}. The geometry and the $3\times3$ matrix
algebra are \emph{the same thing} (Doc.~206); the standard-QM Hilbert space is a \emph{lossless
bijection} (Doc.~230); the information formalism and the spatial reduction are \emph{lossy
projections} with a formalism-relative residue (Doc.~243/270). The decisive point: \emph{only} the
full $T^4$ geometry loses nothing. The θ=2/9 work carries the same result from the other side: all
representations are faithful on the radial (magnitude) sector; the only irreducible part is the
angular (phase) sector.

\section{The question and its boundary}
Over recent months FFGFT has been cast in strikingly many languages, and the impression of several
\enquote{parallel roads} is justified. The question here is internal: \emph{what relationship} holds
between these roads? Are they interchangeable pictures of one object, or do they differ structurally?

This must be set apart from Doc.~265. There a \emph{comparison operator} (three levels --
near-equivalence, structural overlap, complementary; six axes) is developed and applied to FFGFT
\emph{against an external framework} (Blume-Maleev). Here the same operator is turned \emph{inward}:
not FFGFT versus foreign, but FFGFT-representation versus FFGFT-representation. The result is sharper
than the outward operator allows -- and it exposes a relationship type the symmetric outward operator
does not have (Section~\ref{sec:inward}).

\section{The generative core: why $T^4$ is not one representation among many}
Two documents say the same thing from different sides. Doc.~152 (ontological priority) sets the
$T^4$ torsion crystal as \emph{fundamental ontological reality}; the fractal granulation
$D_f=3-\xi$, the effective laws and classical 3D physics are emergent approximations
\enquote{projected upward}. Doc.~243 gives the structural counterpart: every information formalism
discards at least one of the three quantities $\theta_4$, $(n_\theta,n_\varphi)$, $\varepsilon$;
the $T^4$ geometry is the \emph{only} formalism that structurally holds all three at once.

\begin{keyresult}
The $T^4$ torus is not one representation among many but the generative core: the only language that
holds the full content. All other representations are maps \emph{from} $T^4$ -- and split, by the
type of that map, into three classes.
\end{keyresult}

\section{The three relationship types}

\subsection{Identity: geometry $\equiv$ matrix algebra (Doc.~206)}
The Z$_3$ triangle (as a graph) and the $3\times3$ matrix $A$ are not merely isomorphic but
\emph{algebraically identical}: the spectrum of $A$ is the character table of Z$_3$. Doc.~206 puts it
literally: geometry and matrix algebra are \enquote{not two pictures of the same thing, they are the
same thing.} This is the strongest type -- not a lossless translation channel but identity. Honestly
noted: this identity is \emph{sector-local}. It holds on the $C_3$/Z$_3$ core ($3\times3$). Globally
-- with the angular phase $\theta$ and the $\varphi$ harmonics -- it is no longer a finite identity;
the circulant captures the magnitude spectrum, not the phase (Doc.~286).

\subsection{Bijection: the lossless channel (Doc.~230, Doc.~070)}
The Hilbert-space translation (Doc.~230) is a \emph{bijective structural statement}: every QM
calculation sits inside the FFGFT formalism and back, \emph{without loss}. It is no poorer than FFGFT
and no richer in computation -- but FFGFT \emph{adds} what standard QM does not provide: the
\emph{why} of the constant values. In the projection-chain picture (Doc.~270) this is the
\emph{upward} move (Type~III), and precisely because it introduces no geometric extra dimensions it is
lossless.

A second bijection case is corpus-internal: the two mass forms in Doc.~070, the compact (cancelled)
and the extended (geometrically original). They are mathematically equivalent, but only the extended
carries the physical origin -- which is why the ratios must not be cancelled directly.

\subsection{Projection: the lossy, directed channel (Doc.~243, Doc.~270)}
Doc.~243 maps FFGFT into an information formalism -- as a \emph{projection} $\pi$, \emph{not} an
isomorphism. A residue is lost, and that residue is \emph{formalism-relative}: it depends on the
chosen information formalism, not on the geometry. The same mapping into a neural network shares this
character. Geometrically this corresponds to the \emph{spatial} reductions of the projection chain
(Doc.~270, Type~II): strictly lossy, irreversible, with residue $\sim\xi$ (visible as a CHSH
deviation).

To be set apart is the first reduction step $D4\!\to\!D3$ (Doc.~270, Type~I): not a lossy projection
but a \emph{duality decompactification} via $\tilde T\cdot m=1$ -- information-preserving; only the
direct visibility of compactness is lost.

\section{Two axes that are not competing representations}
Not every \enquote{language} is a rival encoding of the same content. Two are orthogonal separations:
\begin{itemize}
  \item \textbf{Two layers} (Doc.~241): the separation of ratio-based, parameter-free content
        (Layer~1) from a single SI anchor (Layer~2). This is bookkeeping -- what is input, what
        follows -- not an alternative geometry. Orthogonal to the geometry/matrix/Hilbert choice.
  \item \textbf{Recursion / fixed point} (Doc.~203/275): the \emph{dynamical} axis. It generates the
        hierarchy (the $\xi\!\to\!\varphi$ passage) and is -- decisively -- purely \emph{real/radial}.
        It carries no angular (phase) generator.
\end{itemize}

\section{The decisive finding}
Two statements that together form the strongest internal-consistency claim.

\begin{important}
\textbf{(1) Only $T^4$ loses nothing.} Every projective encoding drops a residue, and the residue is
relative to the \emph{representation}, not to the physics. The many roads therefore agree; where they
diverge it is exactly what the given encoding threw away. The \enquote{loss} is an artefact of the
representation, not of the theory.

\textbf{(2) The asymmetry (Doc.~270).} Geometrically \emph{downward} is lossy (Type~II);
representationally \emph{upward} to the Hilbert space is lossless (Type~III). This is exactly what
lets FFGFT be \emph{simultaneously} a \enquote{$4$D torus theory} and \enquote{ordinary QM} without
contradiction.
\end{important}

\section{The bracket to the θ=2/9 phase}
The representation analysis and the θ elimination chain (Doc.~282/286) are one result from two sides.
All representations are faithful on the \emph{radial} (magnitude) sector: identity (206) and
bijection (230) capture it completely -- $\xi$, the hierarchy, the Koide amplitude $Q=2/3=r\sqrt2$,
the mean $\langle f\rangle$. The \emph{one} place where the geometry is irreducible is the
\emph{angular} (phase) sector: the offset $\theta=2/9$. Only the full $T^4$ geometry holds it; the
real recursion does not reach it, the circulant does not capture it. Faithful in magnitude,
irreplaceable in phase.

\section{The 265 operator turned inward}\label{sec:inward}
Turning the Doc.~265 comparison operator inward shows two things. First, the internal relationships are
\emph{sharper}: where the outward operator reached only \emph{near}-equivalence (different derivation
routes, identical abstract structure), the internal bijections are \emph{exact} (isomorphism,
Doc.~230) and the matrix case is even \emph{identity} (Doc.~206). Second, the inward application
exposes a type the three symmetric levels of the outward operator do not have: the \emph{directed
projection} (243/270-II). The 265 levels compare two formalisms symmetrically; the internal projection
is directed ($T^4 \to$ encoding) and carries a residue. This is no defect of the 265 operator but a
finding: inward, a fourth, asymmetric relation is needed.

\begin{table}[htbp]
\centering
\caption{The internal representations, by relationship type to the generative $T^4$ core.}
\label{tab:map}
\begin{tabular}{p{4.2cm}p{1.4cm}p{3.0cm}p{4.0cm}}
\toprule
\textbf{Representation} & \textbf{Doc.} & \textbf{Relationship} & \textbf{keeps / loses} \\
\midrule
Z$_3$ triangle $\equiv$ $3\times3$ matrix & 206 & identity & same thing; sector-local ($C_3$) \\
\addlinespace
Hilbert space $H=L^2(T^4)\otimes\mathbb{C}^3$ & 230 & bijection (lossless) & QM $\leftrightarrow$ FFGFT; FFGFT adds the \enquote{why} \\
\addlinespace
compact vs.\ extended mass form & 070 & bijection (equivalent) & equal; only extended carries origin \\
\addlinespace
information formalism & 243 & projection $\pi$ & residue, formalism-relative \\
\addlinespace
neural network & 243 & projection & same \\
\addlinespace
spatial reduction $T^3\!\to\!T^0$ & 270 & projection (Type~II) & residue $\sim\xi$ (CHSH) \\
\midrule
duality decompact.\ $D4\!\to\!D3$ & 270 & Type~I (preserving) & only visibility of compactness \\
\addlinespace
two layers (ratio/SI) & 241 & orthogonal axis & bookkeeping, no rival \\
\addlinespace
recursion / fixed point & 203/275 & orthogonal axis & dynamical, purely real/radial \\
\bottomrule
\end{tabular}
\end{table}

\section{Open edges}
Declared honestly, not hidden:
\begin{itemize}
  \item The \textbf{identity matrix $\equiv$ geometry} (206) is sector-local ($3\times3$, $C_3$). A
        global finite identity including $\theta$ and the $\varphi$ harmonics does not exist; whether
        it \emph{can} exist is open (Doc.~286).
  \item The \textbf{residue table} of the information formalisms (243) is a working stand, not a proof
        of minimality that $T^4$ is the only complete formalism.
  \item The \textbf{265 comparison operator} is itself open with three declared residues
        (axis completeness, weighting, independence); the fourth, directed relation found here is to
        be added.
\end{itemize}

\begin{conclusionBox}[Balance: one core, three relations]
The parallel roads of FFGFT are not of equal rank. The $T^4$ torus is generative -- the only
representation that loses nothing. Everything else stands to it in exactly one of three relations:
\emph{identity} (matrix algebra, sector-local), \emph{bijection} (Hilbert space; compact/extended mass
form), \emph{projection} (information formalism, neural network, spatial reduction -- with a
formalism-relative residue). Two further languages (two layers, recursion) are not rivals but
orthogonal axes. The common ground of all roads is the radial magnitude sector; the only irreducible
place is the angular phase $\theta=2/9$. The representation question is thereby closed as far as it
honestly can be: faithful in magnitude, generative only in the full torus, open only in the phase.
\end{conclusionBox}
